\begin{helper}
Let $F\subset\mathbb{N}$ be finite, let $g_1,g_2\colon\mathbb{N}\to\mathbb{Z}$
and $c_1,c_2\colon\mathbb{N}\to\mathbb{C}$, and assume
\begin{enumerate}
\item[(i)] $g_1$ is injective on $F$: for $k,k'\in F$, $g_1(k)=g_1(k')$ implies
  $k=k'$;
\item[(ii)] $g_2$ is injective on $F$: for $k,k'\in F$, $g_2(k)=g_2(k')$
  implies $k=k'$;
\item[(iii)] the two frequency families never collide: $g_1(k)\neq g_2(k')$ for
  all $k,k'\in F$.
\end{enumerate}
Then
\[
  \int_{0}^{2\pi}\Bigl|\sum_{k\in F}\bigl(c_1(k)\,e^{i g_1(k)\theta}
      +c_2(k)\,e^{i g_2(k)\theta}\bigr)\Bigr|^{2}d\theta
  =2\pi\sum_{k\in F}\bigl(|c_1(k)|^{2}+|c_2(k)|^{2}\bigr).
\]

This is the Parseval step of the energy computation of Section 4
(\texttt{paper.tex:353-357}), in exactly the generality that computation needs.
On the circle of radius $r$ the summand $T_k$ of the series (19) carries
precisely the two frequencies $k-n$ and $-k-n$; the family $k\mapsto k-n$ is
injective, the family $k\mapsto-k-n$ is injective, and the two never meet
because $(k-n)-(-k'-n)=k+k'>0$ for $k,k'\ge n\ge1$.  Applying the identity to
a partial sum of (19) and then integrating in $r$ against the Jacobian is what
produces the $\pi\sum\lambda_k(1+\epsilon_k^{2})|A_k|^{2}$ of the energy
identity.

Hypotheses (i)--(iii) are exactly what makes the diagonal the whole story: the
two injectivity assumptions kill the off-diagonal pairs within each family and
(iii) kills every cross pair.  Nothing is assumed about the sign or the size of
the frequencies, and $F$ is an arbitrary finite index set, so the identity
applies verbatim to any partial sum.
\end{helper}
\begin{proof}
For $l\in\mathbb{Z}$ write $E_l(\theta)=e^{il\theta}$, and abbreviate
\[
  S(\theta)=\sum_{k\in F}\bigl(c_1(k)E_{g_1(k)}(\theta)
    +c_2(k)E_{g_2(k)}(\theta)\bigr).
\]

\emph{Step 0: orthogonality of the characters.}  For every $l\in\mathbb{Z}$,
\begin{equation}\tag{1}
  \int_{0}^{2\pi}E_l(\theta)\,d\theta
  =\begin{cases}2\pi,&l=0,\\[2pt]0,&l\neq0.\end{cases}
\end{equation}
Indeed, if $l=0$ the integrand is the constant $1$ and the integral is $2\pi$.
If $l\neq0$, put $a=il\in\mathbb{C}$, which is nonzero because $l\neq0$ and
$i\neq0$; writing $E_l(\theta)=e^{a\theta}$ and using the antiderivative
$\theta\mapsto e^{a\theta}/a$ of $\theta\mapsto e^{a\theta}$,
\[
  \int_{0}^{2\pi}e^{a\theta}d\theta=\frac{e^{2\pi a}-1}{a}
  =\frac{e^{2\pi i l}-1}{il}=0,
\]
since $e^{2\pi il}=1$ for every integer $l$.

\emph{Step 1: integrability.}  For each $l$ and each $c\in\mathbb{C}$ the
function $\theta\mapsto cE_l(\theta)$ is continuous on $\mathbb{R}$, hence
interval-integrable on $[0,2\pi]$.  Consequently $S$, a finite sum of such
functions, is continuous and interval-integrable there, and so are
$\theta\mapsto\overline{S(\theta)}$ and $\theta\mapsto|S(\theta)|^{2}$.  Every
integral written below is therefore defined, and the integral of a finite sum
of such functions is the sum of their integrals.

\emph{Step 2: expanding the square.}  For every $\theta$ we have
$|S(\theta)|^{2}=S(\theta)\overline{S(\theta)}$, an identity between a real
number and a complex one, read through the inclusion
$\mathbb{R}\hookrightarrow\mathbb{C}$.  Also
\begin{equation}\tag{2}
  \overline{E_l(\theta)}=\overline{e^{il\theta}}=e^{-il\theta}
  =E_{-l}(\theta),
\end{equation}
because conjugation is a continuous field automorphism fixing $\mathbb{R}$ and
sending $i$ to $-i$.  Expanding the product of the two finite sums and using
(2) termwise,
\begin{equation}\tag{3}
  S(\theta)\overline{S(\theta)}
  =\sum_{k\in F}\sum_{k'\in F}
    \begin{aligned}[t]
      \Bigl[&c_1(k)\overline{c_1(k')}\,E_{g_1(k)-g_1(k')}(\theta)
       +c_1(k)\overline{c_2(k')}\,E_{g_1(k)-g_2(k')}(\theta)\\
      &{}+c_2(k)\overline{c_1(k')}\,E_{g_2(k)-g_1(k')}(\theta)
       +c_2(k)\overline{c_2(k')}\,E_{g_2(k)-g_2(k')}(\theta)\Bigr],
    \end{aligned}
\end{equation}
using $E_a(\theta)E_b(\theta)=E_{a+b}(\theta)$, which is the additivity of the
exponential.

\emph{Step 3: integrating, and discarding the off-diagonal terms.}  Integrate
(3) over $[0,2\pi]$.  The right-hand side is a finite sum of functions of the
form $cE_l$, so by Step 1 the integral is the corresponding finite sum of
integrals, and each is $c\int_0^{2\pi}E_l$, evaluated by (1).  We examine the
four families of terms.

The first family contributes only when $g_1(k)-g_1(k')=0$, i.e.
$g_1(k)=g_1(k')$, which by (i) happens exactly when $k=k'$; for such a pair
the contribution is $2\pi\,c_1(k)\overline{c_1(k)}=2\pi|c_1(k)|^{2}$.

The fourth family likewise contributes only when $g_2(k)=g_2(k')$, which by
(ii) happens exactly when $k=k'$, with contribution
$2\pi\,c_2(k)\overline{c_2(k)}=2\pi|c_2(k)|^{2}$.

The second family would need $g_1(k)=g_2(k')$ for some $k,k'\in F$, which (iii)
forbids, so every one of its terms integrates to $0$.  The third family would
need $g_2(k)=g_1(k')$ for some $k,k'\in F$; applying (iii) with the two
indices interchanged gives $g_1(k')\neq g_2(k)$, so this too is impossible and
every term integrates to $0$.

Hence only the diagonal pairs $k=k'$ of the first and fourth families survive,
and
\begin{equation}\tag{4}
  \int_{0}^{2\pi}S(\theta)\overline{S(\theta)}\,d\theta
  =\sum_{k\in F}2\pi\bigl(|c_1(k)|^{2}+|c_2(k)|^{2}\bigr)
  =2\pi\sum_{k\in F}\bigl(|c_1(k)|^{2}+|c_2(k)|^{2}\bigr).
\end{equation}

\emph{Step 4: descending to the real identity.}  By Step 2 the integrand on
the left of (4) is the image of the real-valued function
$\theta\mapsto|S(\theta)|^{2}$ under the inclusion
$\mathbb{R}\hookrightarrow\mathbb{C}$, and that inclusion commutes with the
integral of an interval-integrable real function.  So the left side of (4) is
the image of $\int_{0}^{2\pi}|S(\theta)|^{2}d\theta$, while the right side is
the image of the real number
$2\pi\sum_{k\in F}(|c_1(k)|^{2}+|c_2(k)|^{2})$.  The inclusion
$\mathbb{R}\hookrightarrow\mathbb{C}$ is injective, so the two real numbers are
equal, which is the assertion.
\end{proof}
